\documentclass[10pt,a4paper]{amsart}
\usepackage[margin=19mm]{geometry}
\usepackage{amsmath,amssymb,mathtools}
\usepackage[scr=rsfs,cal=euler]{mathalfa}
\usepackage{xfrac}
\usepackage[unicode,hidelinks,bookmarks=false]{hyperref}
\newcommand{\ie}{i.e.\ }
\newcommand{\cf}{cf.\ }
\newcommand{\resp}{respectively\ }
\newcommand{\Subsection}{\subsection}
\providecommand{\nobreakdash}{\nobreak\hbox{-}}
\setlength{\emergencystretch}{2em}
\multlinegap0pt
\usepackage{bm}
\usepackage{bbm}
\usepackage{dsfont}
\usepackage{xspace}
\usepackage{cancel}
\usepackage{mathtools}
\usepackage{amsthm}
\makeatletter
\@ifundefined{theorem}{\newtheorem{theorem}{Theorem}}{}
\@ifundefined{lemma}{\newtheorem{lemma}[theorem]{Lemma}}{}
\makeatother
\newcommand{\CC}{\mathbb{C}}
\title[On Extremal Family Trees $(\mathcal{T}\_n)\_{n\geqslant 3}$ Be]{On Extremal Family Trees $(\mathcal{T}\_n)\_{n\geqslant 3}$ Beyond Caterpillars and Greedy Constructions}
\author{Jasem Hamoud, Duaa Abdullah}
\date{Source: arXiv:2602.03909v1 (CC BY 4.0)}
\begin{document}
\maketitle

\noindent\emph{Source and licence.} On Extremal Family Trees $(\mathcal{T}_n)_{n\geqslant 3}$ Beyond Caterpillars and Greedy Constructions. Version arXiv:2602.03909v1 (CC BY 4.0), distributed under the Creative Commons Attribution 4.0 International licence, as recorded in the arXiv licence metadata for this exact version and shown on the abstract page https://arxiv.org/abs/2602.03909v1. Full rights notice: licenses/arxiv-2602.03909-CC-BY-4.0.txt.\par\smallskip

\noindent\textit{Editorial scope.} Counted unit: the Theorem whose source label is \texttt{le.sigma2} in the article (source0.tex lines 583--591, source label \texttt{le.sigma2}), delivered together with the article's own complete original proof of it (source0.tex lines 593--624, one \texttt{proof} environment of 32 lines). No mathematics is written, repaired or re-derived: statement and proof are the article's own text line for line. Required context retained uncounted: le.segma3 (lines 458--463). Every definition, display and label of those lines is retained unchanged. Omitted from this package and not redistributed: 1--457, 464--582, 592--592, 625--871. Nothing inside a retained excerpt is omitted or altered: every retained body is delivered exactly as the article's lines are written, so no omission receipt of any kind accompanies this package. Citations: the retained excerpts carry no citation command at all, so no bibliography entry of the article is transcribed and no reference is redistributed. Reference adaptation: none was needed and none was made; every reference inside the delivered unit resolves inside it, so no unresolved reference or citation is produced. Printed slips kept literal: no printed word, symbol, hypothesis or number of the counted statement or of its proof was altered. Layout and class: the article's own class is \texttt{article} with packages including graphicx, amssymb, stmaryrd, enumitem, geometry, tikz, float, url, mathtools, amsmath, amsfonts, amsthm, mathrsfs, pgfplots; the wrapper is the lane's pinned \texttt{amsart} editorial wrapper at 10pt on A4, which restates the theorem-like environments the retained text needs and adds the article's own macro definitions that the retained text needs (\texttt{\textbackslash CC}), with extra wrapper packages (bm, bbm, dsfont, xspace, cancel, mathtools). The delivered numbering is the unit's own. Assets: no retained span contains a figure environment, a graphics-inclusion command, a PDF-page inclusion, a table or a TikZ picture; no image file is copied into this package, the package declares no asset dependency and no image file is redistributed with it. Licence: version arXiv:2602.03909v1 of this article is distributed under the Creative Commons Attribution 4.0 International licence, as recorded in the arXiv licence metadata for this exact version and shown on the abstract page https://arxiv.org/abs/2602.03909v1. A full rights notice is delivered at licenses/arxiv-2602.03909-CC-BY-4.0.txt. Characters: every retained excerpt is pure ASCII, so the delivered engine, XeLaTeX with its default Latin Modern text font, reports no missing character.
\begin{lemma}\label{le.segma3}
Let $\CC$ be a caterpillar tree with degree sequence $ \mathscr{D} = (d_1, d_2, d_3) $, where $ d_3 \geqslant d_2 \geqslant d_1 $. Then, the Sigma index of $ T $ is given by:
\begin{equation}~\label{eqle.segma3n1}
\sigma(\CC)=\sum_{i=1}^{3} p\,(d_i-1)^2+\sum_{i=1}^{2}(d_i-d_{i+1})^2.
\end{equation}
\end{lemma}

\begin{theorem}~\label{le.sigma2}
Let $\CC$ be a caterpillar tree of order $n$, and let $\mathscr{D} = (d_1, d_2, d_3, d_4)$ be a degree sequence, where $d_4 \geqslant d_3 \geqslant d_2 \geqslant d_1$. Then, 
\[
\sigma(\CC)=\begin{cases}
\sigma_{\max}(\CC)=p \sum_{i=1}^{4} (d_i-1)^2+(d_4-d_1)^2+(d_1-d_3)^2+(d_3-d_2)^2, \\
 \sigma_{\min}(\CC)= p \sum_{i=1}^{4} (d_i-1)^2+(d_1-d_2)^2+(d_2-d_3)^2+(d_3-d_4)^2.
\end{cases}
\]
\end{theorem}

\begin{proof}
Let $\CC$ be a caterpillar tree of order $n$, and let $\mathscr{D} = (d_1, d_2, d_3, d_4)$ be a non-increasing degree sequence. We aim to determine which vertices at the ends of the sequence represent the largest values, and which vertices in the middle represent the smallest values.  Assume the degrees satisfy $a \leq b \leq c \leq d$, where $c$ and $d$ are the two largest. For any ordering $(x_1,x_2,x_3,x_4)$, the Sigma index is
$$
\sigma(\CC) = p\sum_{i=1}^{4}(x_i-1)^2 + \sum_{i=1}^{3}(x_i-x_{i+1})^2.
$$
The term $p\sum_{i=1}^{4}(x_i-1)^2$ depends only on the multiset $\{a,b,c,d\}$ and remains constant across all permutations. The condition that the two largest degrees are adjacent requires $(c,d)$ or $(d,c)$ to appear consecutively in $(x_1,x_2,x_3,x_4)$. Thus, we noticed that The six corresponding orderings are
$(c,d,a,b),\ (d,c,a,b),\ (a,c,d,b),$ and $(a,d,c,b),\ (a,b,c,d),\ (a,b,d,c).$ For each ordering, according to Lemma~\ref{le.segma3} we noticed that
\begin{equation}~\label{eqq1le.sigma2}
\sigma(\CC)= p\big[(a-1)^2+(b-1)^2+(c-1)^2+(d-1)^2\big] + \mathcal{S},
\end{equation}

where $\mathcal{S}$ denotes the sum of squared consecutive differences:
\begin{align*}
 &\mathcal{S}_{(a,b,c,d)}= (a-b)^2 + (b-c)^2 + (c-d)^2, \quad 
\mathcal{S}_{(a,b,d,c)}= (a-b)^2 + (b-d)^2 + (d-c)^2,\\
&\mathcal{S}_{(a,c,d,b)}= (a-c)^2 + (c-d)^2 + (d-b)^2,\quad
\mathcal{S}_{(a,d,c,b)}= (a-d)^2 + (d-c)^2 + (c-b)^2,\\
&\mathcal{S}_{(c,d,a,b)}= (c-d)^2 + (d-a)^2 + (a-b)^2,\quad
\mathcal{S}_{(d,c,a,b)}= (d-c)^2 + (c-a)^2 + (a-b)^2.
\end{align*}
In each case, the term involving the two largest degrees is either $(c-d)^2$ or $(d-c)^2$, appearing exactly once since $c$ and $d$ are adjacent only once in a sequence of length $4$.
Thus, the maximum value of Sigma index according to~\eqref{eqq1le.sigma2} is 
\begin{equation}
\sigma_{\max}(\CC)=p \sum_{i=1}^{4} (x_i - 1)^2 + (d-a)^2 + (a-c)^2 + (c-b)^2.
\end{equation}
Adjacent terms are closest, minimizing the sum of squared differences. Therefore, the minimum value  is 
\begin{equation}~\label{eqq3le.sigma2}
 \sigma_{\min}(\CC)= p \sum_{i=1}^{4} (x_i - 1)^2 + (a-b)^2 + (b-c)^2 + (c-d)^2.
\end{equation}
Thus, from~\eqref{eqq1le.sigma2}--\eqref{eqq3le.sigma2} holds.
This completes the proof.
\end{proof}

\end{document}
